\section{AMC 10 10A 2005}
        \begin{enumerate}
        	\item (\textit{AMC 10 10A 2005}) While eating out, Mike and Joe each tipped their server $\$2$. Mike tipped $10\%$ of his bill and Joe tipped $20\%$ of his bill. What was the difference, in dollars, between their bills? 


 $\textbf{(A) } 2\qquad \textbf{(B) } 4\qquad \textbf{(C) } 5\qquad \textbf{(D) } 10\qquad \textbf{(E) } 20$


 
 \textbf{Answer: }D



\textbf{Solution}

Let $m$ be Mike's bill and $j$ be Joe's bill (both in dollars). We then have $\frac{10}{100}m = 2 \iff m = 20$ and $\frac{20}{100}j = 2 \iff j = 10$, so the desired difference is $\left\lvert m-j\right\rvert = 20-10 = \boxed{\textbf{(D) } 10}$.

 


	\item (\textit{AMC 10 10A 2005}) For each pair of real numbers $a \neq b$, define the operation $\star$ as


 
 \[(a \star b) = \frac{a+b}{a-b}.\]


 What is the value of $\left(\left(1 \star 2\right) \star 3\right)$?


 $\textbf{(A) } -\frac{2}{3}\qquad \textbf{(B) } -\frac{1}{5}\qquad \textbf{(C) } 0\qquad \textbf{(D) } \frac{1}{2}\qquad \textbf{(E) } \text{This value is not defined.}$


 
 \textbf{Answer: }C



\textbf{Solution}

$((1 \star 2) \star 3) = \left(\left(\frac{1+2}{1-2}\right) \star 3\right) = (-3 \star 3) = \frac{-3+3}{-3-3} = \boxed{\textbf{(C) } 0}$.

 


	\item (\textit{AMC 10 10A 2005}) The equations $2x + 7 = 3$ and $bx - 10 = -2$ have the same solution $x$. What is the value of $b$? 


 $\textbf{(A) } -8\qquad \textbf{(B) } -4\qquad \textbf{(C) } -2\qquad \textbf{(D) } 4\qquad \textbf{(E) } 8$


 
 \textbf{Answer: }B



\textbf{Solution}

$2x + 7 = 3 \iff x = -2$, so we require $-2b-10 = -2 \iff -2b = 8 \iff b = \boxed{\textbf{(B) } -4}$.

 


	\item (\textit{AMC 10 10A 2005}) A rectangle with a diagonal of length $x$ is twice as long as it is wide. What is the area of the rectangle? 


 $\textbf{(A) } \frac{1}{4}x^2\qquad \textbf{(B) } \frac{2}{5}x^2\qquad \textbf{(C) } \frac{1}{2}x^2\qquad \textbf{(D) } x^2\qquad \textbf{(E) } \frac{3}{2}x^2$


 
 \textbf{Answer: }B



\textbf{Solution 1}

Let's set the length to $2$ and the width to $1$, so the rectangle has area $2 \cdot 1 = 2$ and diagonal $x = \sqrt{1^2+2^2} = \sqrt{5}$ (by Pythagoras' theorem).

 Now we can plug this value into the answer choices and test which one will give our desired area of $2$. (Note that all of the answer choices contain $x^2$, so keep in mind that $\left(\sqrt{5}\right)^2 = 5$.)

 We eventually find that $\frac{2}{5} \cdot \left(\sqrt{5}\right)^2 = 2$, so the correct answer is $\boxed{\textbf{(B) } \frac{2}{5}x^2}$.

 -JinhoK

 



\textbf{Solution 2}

Call the length $2l$ and the width $l$, so the area of the rectangle is $2l \cdot l = 2l^2$. As in Solution 1, $x$ is the hypotenuse of the right triangle with legs $2l$ and $l$, so by Pythagoras' theorem, $x = \sqrt{\left(2l\right)^2+l^2} = \sqrt{4l^2+l^2} = \sqrt{5l^2}$. This becomes $l^2 = \frac{x^2}{5}$, and therefore the area is $2l^2 = \boxed{\textbf{(B) } \frac{2}{5}x^2}$.

 -mobius247

 


	\item (\textit{AMC 10 10A 2005}) A store normally sells windows at $\$100$ each. This week the store is offering one free window for each purchase of four. Dave needs seven windows and Doug needs eight windows. How many dollars will they save if they purchase the windows together rather than separately?


 $\textbf{(A) } 100\qquad \textbf{(B) } 200\qquad \textbf{(C) } 300\qquad \textbf{(D) } 400\qquad \textbf{(E) } 500$


 
 \textbf{Answer: }A



\textbf{Solution}

The store's offer means that every $5$th window is free, so Dave would get $\left\lfloor\frac{7}{5}\right\rfloor=1$ free window and Doug would also get $\left\lfloor\frac{8}{5}\right\rfloor=1$ free window. 

 This is a total of $1+1 = 2$ free windows. 

 Together, by contrast, they would get $\left\lfloor\frac{8+7}{5}\right\rfloor = \left\lfloor\frac{15}{5}\right\rfloor=3$ free windows.

 Hence they get $3-2=1$ additional window if they purchase the windows together, so the amount (in dollars) that they save is $1\cdot100= \boxed{\textbf{(A) } 100}$.

 


	\item (\textit{AMC 10 10A 2005}) The average (mean) of $20$ numbers is $30$, and the average of $30$ other numbers is $20$. What is the average of all $50$ numbers?


 $\textbf{(A) } 23\qquad \textbf{(B) } 24\qquad \textbf{(C) } 25\qquad \textbf{(D) } 26\qquad \textbf{(E) } 27$


 
 \textbf{Answer: }B



\textbf{Solution}

Since the average of the first $20$ numbers is $30$, their sum is $20\cdot30=600$. 

 Since the average of $30$ other numbers is $20$, their sum is $30\cdot20=600$. 

 So the sum of all $50$ numbers is $600+600=1200$.

 Therefore, the average of all $50$ numbers is $\frac{1200}{50}=\boxed{\textbf{(B) } 24}$.

 


	\item (\textit{AMC 10 10A 2005}) Josh and Mike live $13$ miles apart. Josh started to ride his bicycle toward Mike's house. A little later Mike started to ride his bicycle toward Josh's house. When they met, Josh had ridden for twice the length of time as Mike and at four-fifths of Mike's rate. How many miles had Mike ridden when they met? 


 $\textbf{(A) } 4\qquad \textbf{(B) } 5\qquad \textbf{(C) } 6\qquad \textbf{(D) } 7\qquad \textbf{(E) } 8$


 
 \textbf{Answer: }B



\textbf{Solution}

Let $m$ be the distance in miles that Mike rode. 

 Since Josh rode for twice the length of time as Mike and at four-fifths of Mike's rate, he rode $2\cdot\frac{4}{5}\cdot m = \frac{8}{5}m$ miles. 

 Now, as their combined distance was $13$ miles, we deduce that $\frac{8}{5}m + m = 13 \iff \frac{13}{5}m = 13 \iff m = \boxed{\textbf{(B) }5}$.

 


	\item (\textit{AMC 10 10A 2005}) In the figure, the length of side $AB$ of square $ABCD$ is $\sqrt{50}$, $E$ is between $B$ and $H$, and $BE = 1$. What is the area of the inner square $EFGH$?


 
                    \begin{figure}[H]
                        \centering
                        \includegraphics[width=0.5\textwidth]{Problem/AMC_10_10A_2005_8_Problem_Figure_1.png}
                    \end{figure}
                    


 $\textbf{(A) } 25\qquad\textbf{(B) } 32\qquad\textbf{(C) } 36\qquad\textbf{(D) } 40\qquad\textbf{(E) } 42$


 
 \textbf{Answer: }C



\textbf{Solution}

We see that side $BE$, which we know is $1$, is also the shorter leg of one of the four right triangles (which are obviously congruent, using the symmetry of the diagram). So $AH = 1$, and hence $HB = HE + BE = HE + 1$. Since $HE$ is one of the sides of the square whose area we want to find, we can now simply apply Pythagoras' theorem:

 
 \begin{align*}1^2 + (HE+1)^2 = \left(\sqrt{50}\right)^2 &\iff 1 + (HE+1)^2 = 50 \\ &\iff (HE+1)^2 = 49 \\&\iff HE+1 = 7 \qquad \text{(as } HE \text{ is a length, so must be positive)} \\ &\iff HE = 6.\end{align*}
 

 Thus the area of the square is $6^2 = \boxed{\textbf{(C) }36}$.

 


	\item (\textit{AMC 10 10A 2005}) Three tiles are marked $X$ and two other tiles are marked $O$. The five tiles are randomly arranged in a row. What is the probability that the arrangement reads $XOXOX$?


 $\textbf{(A) } \frac{1}{12}\qquad \textbf{(B) } \frac{1}{10}\qquad \textbf{(C) } \frac{1}{6}\qquad \textbf{(D) } \frac{1}{4}\qquad \textbf{(E) } \frac{1}{3}$


 
 \textbf{Answer: }B



\textbf{Solution 1}

There are $\frac{5!}{2!3!} = 10$ distinct arrangements of $3$ $X$s and $2$ $O$s, and only $1$ distinct arrangement that reads $XOXOX$.

 Therefore the desired probability is simply $\boxed{\textbf{(B) }\frac{1}{10}}$.

 



\textbf{Solution 2}

Arraging the $3$ $X$s and $2$ $O$s in a row is equivalent to fitting the $O$s into the gaps between the $X$s.

 There are a total of $3+1 = 4$ gaps between the $X$s, so fitting $2$ $O$s into these gaps gives $2$ possible outcomes:

 1. The $2$ $O$s are put into different gaps. In this case, the number of possible arrangements is $\binom{4}{2} = 6$.

 2. The $2$ $O$s are put into the same gap. In this case, as there are $4$ gaps, we simply obtain $4$ possible arrangements.

 Therefore the probability of an arrangement that reads $XOXOX$ is $\frac{1}{4 + 6} = \boxed{\textbf{(B) }\frac{1}{10}}$.

 


	\item (\textit{AMC 10 10A 2005}) There are two values of $a$ for which the equation $4x^2 + ax + 8x + 9 = 0$ has only one solution for $x$. What is the sum of those values of $a$?


 $\textbf{(A) } -16\qquad \textbf{(B) } -8\qquad \textbf{(C) } 0\qquad \textbf{(D) } 8\qquad \textbf{(E) } 20$


 
 \textbf{Answer: }A



\textbf{Solution 1}

A quadratic equation has exactly $1$ distinct root if and only if the left-hand side is a perfect square. So we require

 
 \[4x^2 + ax + 8x + 9 = (mx + n)^2 = m^2 x^2 + 2mnx + n^2.\]

 Two polynomials are equal only if their coefficients are equal, so we must have $m^2 = 4$ and $n^2 = 9$, which reduce to $m = \pm 2$ and $n = \pm 3$ respectively. By equating $x$-coefficients, it follows that $a + 8 = 2mn = \pm 2 \cdot 2\cdot 3 = \pm 12$, so either $a = 12-8 = 4$ or $a = -12-8 = -20$.

 Accordingly, the desired sum is $4+(-20) = \boxed{\textbf{(A) } -16}$.

 (Alternatively, we can observe that whatever the $2$ values of $a$ are, they must lead to equations of the form $px^2 + qx + r = 0$ and $px^2 - qx + r = 0$, corresponding to the perfect squares $\left(mx \pm n\right)^2 = 0$. So the $2$ choices of $a$, say $a_1$ and $a_2$, must satisfy $a_1 + 8 = q$ and $a_2 + 8 = -q$, and adding these equations yields $a_1 + a_2 + 16 = 0 \iff a_1 + a_2 = \boxed{\textbf{(A) } -16}$.)

 



\textbf{Solution 2}

Since this quadratic must have a double root, its discriminant must be $0$. Therefore we require
 \[(a+8)^2 - 4 \cdot 4 \cdot  9 = 0 \iff a^2 + 16a - 80 = 0.\] At this point, Vieta's formulae directly give the sum of the $2$ possible values of $a$ as $-\frac{16}{1} = \boxed{\textbf{(A) } -16}.$

 (Alternatively, we could use the quadratic formula to directly solve this equation for $a$; the precise value of the expression under the square root does not actually matter, as it will cancel out due to the $\pm$ signs when added. This again gives the required sum as
 \[\frac{-16 + \sqrt{\text{something}}}{2} + \frac{-16 - \sqrt{\text{something}}}{2} = \frac{-32}{2} = \boxed{\textbf{(A) } -16}.)\]

 


	\item (\textit{AMC 10 10A 2005}) A wooden cube $n$ units on a side is painted red on all six faces and then cut into $n^3$ unit cubes. Exactly one-fourth of the total number of faces of the unit cubes are red. What is $n$?


 $\textbf{(A) } 3\qquad \textbf{(B) } 4\qquad \textbf{(C) } 5\qquad \textbf{(D) } 6\qquad \textbf{(E) } 7$


 
 \textbf{Answer: }B



\textbf{Solution 1}

Since there are $n^2$ little faces on each face of the big wooden cube, there are a total of $6n^2$ little faces painted red. Moreover, as each unit cube has $6$ faces, there are $6n^3$ little faces in total. 

 Accordingly, as one-fourth of the little faces are painted red, we have
 \[\frac{6n^2}{6n^3} = \frac{1}{4} \iff \frac{1}{n} = \frac{1}{4} \iff n = \boxed{\textbf{(B) } 4}.\]

 



\textbf{Solution 2}

We recall that when a cube of side length $n$ has its entire surface painted and is then split into $n^3$ unit cubes, there will be exactly $(n-2)^3$ unit cubes with $0$ faces painted, $6(n-2)^2$ unit cubes with $1$ face painted, $12(n-2)$ unit cubes with $2$ faces painted, and $8$ unit cubes with $3$ faces painted.

 (Observe that these form the terms of the binomial expansion of $n^3 = \left(\left(n-2\right)+2\right)^3$.)

 Hence the total number of faces painted red is 
 \begin{align*}(n-2)^3 \cdot 0 + 6(n-2)^2 \cdot 1 + 12(n-2) \cdot 2 + 8 \cdot 3 &= 6\left(n^2-4n+4\right)+24n-48+24 \\ &= 6n^2-24n+24+24n-24 \\ &= 6n^2,\end{align*}
 and now we may proceed as in Solution 1, again giving $n = \boxed{\textbf{(B) } 4}.$

 


	\item (\textit{AMC 10 10A 2005}) The figure shown is called a \textit{trefoil} and is constructed by drawing circular sectors about sides of the congruent equilateral triangles. What is the area of a trefoil whose horizontal base has length $2$?


 
                    \begin{figure}[H]
                        \centering
                        \includegraphics[width=0.5\textwidth]{Problem/AMC_10_10A_2005_12_Problem_Figure_1.png}
                    \end{figure}
                    


 $\textbf{(A) } \frac{1}{3}\pi+\frac{\sqrt{3}}{2}\qquad \textbf{(B) } \frac{2}{3}\pi\qquad \textbf{(C) } \frac{2}{3}\pi+\frac{\sqrt{3}}{4}\qquad \textbf{(D) } \frac{2}{3}\pi+\frac{\sqrt{3}}{3}\qquad \textbf{(E) } \frac{2}{3}\pi+\frac{\sqrt{3}}{2}$


 
 \textbf{Answer: }B



\textbf{Solution}

The area of the \textit{trefoil} is equal to the area of an equilateral triangle with side length $2$, plus the area of $4$ segments. Each segment has area equal to that of a $60^{\circ}$ sector with radius $\frac{2}{2} = 1$, minus the area of an equilateral triangle with side length $1$.

 As there are $4$ segments, the area of the equilateral triangle with side length $1$ will be multiplied by $4$, and this is equivalent to the area of an equilateral triangle with side length $1 \cdot \sqrt{4} = 2$ (since the area scale factor is the square of the length scale factor). Accordingly, the area of the equilateral triangle with side length $2$ will exactly cancel out, and we are left only with $4$ times the area of a $60^{\circ}$ sector with radius $1$. 

 Thus the answer is $4\cdot\frac{60}{360}\cdot\pi\cdot 1^2 = \frac{4}{6}\cdot\pi = \boxed{\textbf{(B) } \frac{2}{3}\pi}$.

 


	\item (\textit{AMC 10 10A 2005}) How many positive integers $n$ satisfy the following condition:


 
 \[\left(130n\right)^{50} > n^{100} > 2^{200} \ \text{?}\]


 $\textbf{(A) } 0\qquad \textbf{(B) } 7\qquad \textbf{(C) } 12\qquad \textbf{(D) } 65\qquad \textbf{(E) } 125$


 
 \textbf{Answer: }E



\textbf{Solution 1}

Since $n > 0$, all $3$ terms of the inequality are positive, so we may take the $50$th root, yielding

 
 \begin{align*}&\sqrt[50]{(130n)^{50}} > \sqrt[50]{n^{100}} > \sqrt[50]{2^{200}} \\ \iff &130n > n^2 > 2^4 \\ \iff &130n > n^2 > 16.\end{align*}


 Solving each part separately, while noting that $n > 0$, therefore gives $n^2 > 16 \iff n > 4$ and $130n > n^2 \iff n < 130$. 

 Hence the solution is $4 < n < 130$, and therefore the answer is the number of positive integers in the open interval $(4,130)$, which is $129-5+1 = \boxed{\textbf{(E) } 125}$.

 


	\item (\textit{AMC 10 10A 2005}) How many three-digit numbers satisfy the property that the middle digit is the average of the first and the last digits? 


 $\textbf{(A) } 41\qquad \textbf{(B) } 42\qquad \textbf{(C) } 43\qquad \textbf{(D) } 44\qquad \textbf{(E) } 45$


 
 \textbf{Answer: }E



\textbf{Solution 1}

If the middle digit is the average of the first and last digits, twice the middle digit must be equal to the sum of the first and last digits. 

 Doing some \href{https://artofproblemsolving.com/wiki/index.php?title=Casework}{casework}: 

 If the middle digit is $1$, possible numbers range from $111$ to $210$. So there are $2$ numbers in this case. 

 If the middle digit is $2$, possible numbers range from $123$ to $420$. So there are $4$ numbers in this case. 

 If the middle digit is $3$, possible numbers range from $135$ to $630$. So there are $6$ numbers in this case. 

 If the middle digit is $4$, possible numbers range from $147$ to $840$. So there are $8$ numbers in this case. 

 If the middle digit is $5$, possible numbers range from $159$ to $951$. So there are $9$ numbers in this case. 

 If the middle digit is $6$, possible numbers range from $369$ to $963$. So there are $7$ numbers in this case. 

 If the middle digit is $7$, possible numbers range from $579$ to $975$. So there are $5$ numbers in this case. 

 If the middle digit is $8$, possible numbers range from $789$ to $987$. So there are $3$ numbers in this case. 

 If the middle digit is $9$, the only possible number is $999$. So there is $1$ number in this case. 

 So the total number of three-digit numbers that satisfy the property is $2+4+6+8+9+7+5+3+1=\boxed{\textbf{(E) } 45}$.

 



\textbf{Solution 2}

Alternatively, we could note that the middle digit is uniquely defined by the first and third digits since it is half of their sum. This also means that the sum of the first and third digits must be even. Since even numbers are formed either by adding two odd numbers or two even numbers, we can split our problem into $2$ cases:

 If both the first digit and the last digit are odd, then we have $1, 3, 5, 7,$ or $9$ as choices for each of these digits, and there are $5\cdot5=25$ numbers in this case.

 If both the first and last digits are even, then we have $2, 4, 6, 8$ as our choices for the first digit and $0, 2, 4, 6, 8$ for the third digit. There are $4\cdot5=20$ numbers here.

 The total number, then, is $20+25=\boxed{\textbf{(E) } 45}$.

 



\textbf{Solution 3}

As we noted in Solution 2, we note that the sum of the first and third digits has to be even. The first digit can have $9$ possibilities $(1-9)$, and the third digit can have $10$ possibilities $(0-9)$. This means there can be $9\cdot10=90$ possible two-digit numbers in which the first digit and the third digit are digits. Exactly half of these would have their sum be divisible by $2$ (since $90$ is even), so our answer is $\frac{90}{2}=\boxed{\textbf{(E) }45}$.

 - \href{https://artofproblemsolving.com/wiki/index.php?title=User:Sweetmango77}{SweetMango77}

 



\textbf{Solution 4}

If our first digit is odd, the last digit also has to be odd for there to be a whole number middle digit that is the average of the first and the last digits.

 If our first digit is even, the last digit has to be even, or $0$ for there to be a whole number middle digit that is the average of the first and the last digits.

 As there are $5$ different possible values for each starting digit of our number, we can multiply that by the number of possible first digits $(1-9)$ to achieve our desired solution. $5 \cdot 9 = \boxed{\textbf{(E) }45}$.

 -JinhoK

 


	\item (\textit{AMC 10 10A 2005}) How many positive cubes divide $3! \cdot 5! \cdot 7!$?


 $\textbf{(A) } 2\qquad \textbf{(B) } 3\qquad \textbf{(C) } 4\qquad \textbf{(D) } 5\qquad \textbf{(E) } 6$


 
 \textbf{Answer: }E



\textbf{Solution 1}

We obtain $3! \cdot 5! \cdot 7! = (3\cdot2\cdot1) \cdot (5\cdot4\cdot3\cdot2\cdot1) \cdot (7\cdot6\cdot5\cdot4\cdot3\cdot2\cdot1) = 2^{8}\cdot3^{4}\cdot5^{2}\cdot7$.

 Therefore, a perfect cube that divides $3! \cdot 5! \cdot 7!$ must be of the form $2^{a}\cdot3^{b}\cdot5^{c}\cdot7^{d}$, where $a$, $b$, $c$, and $d$ are non-negative multiples of $3$ that are less than or equal to $8$, $4$, $2$ and $1$ respectively. 

 This means 

 
 \begin{align*}&a\in\{0,3,6\} \text{ (} 3 \text{ possibilities),} \\ &b\in\{0,3\} \text{ (} 2 \text{ possibilities),} \\ &c\in\{0\} \text{ (} 1 \text{ possibility), and} \\ &d\in\{0\} \text{ (} 1 \text{ possibility),}\end{align*}


 so the number of perfect cubes that divide $3! \cdot 5! \cdot 7!$ is $3\cdot2\cdot1\cdot1 = \boxed{\textbf{(E) } 6}$.

 



\textbf{Solution 2}

As in Solution 1, we write $3! \cdot 5! \cdot 7! = (3\cdot2\cdot1) \cdot (5\cdot4\cdot3\cdot2\cdot1) \cdot (7\cdot6\cdot5\cdot4\cdot3\cdot2\cdot1)$, and now notice that there are a total of $3$ $3$s, $3$ $2$s, and $3$ $1$s. This gives us our first $3$ cubes: $3^3$, $2^3$, and $1^3$. 

 However, we can also multiply smaller numbers in the expression to make bigger cubes. For example, $(2 \cdot 2) \cdot 4 \cdot 4 = 4 \cdot 4 \cdot 4 = 4^3$ (where one 2 comes from the $3!$, and the other from the $5!$). Using this method, we can also make

 
 \[(3 \cdot 2) \cdot (3 \cdot 2) \cdot 6 = 6^3\]

 and

 
 \[(3 \cdot 4) \cdot (3 \cdot 4) \cdot (2 \cdot 6) = 12^3.\]

 So we have $6$ possible cubes in total: $1^3$, $2^3$, $3^3$, $4^3$, $6^3$, and $12^3$, and the answer is $\boxed{\textbf{(E) } 6}$.

 


	\item (\textit{AMC 10 10A 2005}) The sum of the digits of a two-digit number is subtracted from the number. The units digit of the result is $6$. How many two-digit numbers have this property? 


 $\textbf{(A) } 5\qquad \textbf{(B) } 7\qquad \textbf{(C) } 9\qquad \textbf{(D) } 10\qquad \textbf{(E) } 19$


 
 \textbf{Answer: }D



\textbf{Solution 1}

Let the number be $10a+b$, where $a$ is its tens digit and $b$ is its units digit. Then $(10a+b)-(a+b) = 9a$ must have a units digit of $6$, and as $a$ is the tens digit, we can only have $1 \leq a \leq 9$, so $9 \leq 9a \leq 81$.

 Accordingly, $9a$ has units digit $6$ only if $9a = 36 \iff a=4$. Thus the numbers that have the required property are all those with tens digit $4$, from $40$ to $49$, so the answer is $49-40+1 = \boxed{\textbf{(D) } 10}$.

 



\textbf{Solution 2}

As in Solution 1, suppose that $a$ and $b$ are the tens and units digits of the number respectively, so the result of the subtraction is $10a+b-(a+b)=9a$. Thus $9a$ must have units digit $6$.

 We now observe that $10a$ is a multiple of 10, so has units digit $0$, and hence $10a-9a = a$ will have units digit $10-6 = 4$ (where the $0$ will become $10$ in the subtraction by 'borrowing' $1$ from the tens digit). Since $a$ is a single digit, this simply means $a = 4$, while $b$ can be any digit from $0$ to $9$ (since it cancelled out in the subtraction above).

 Thus, as in Solution 1, there are $10$ possible choices for $b$ with $a = 4$, so the answer is $\boxed{\textbf{(D) } 10}$.

 


	\item (\textit{AMC 10 10A 2005}) In the five-sided star shown, the letters $A$, $B$, $C$, $D$, and $E$ are replaced by the numbers $3$, $5$, $6$, $7$, and $9$, although not necessarily in this order. The sums of the numbers at the ends of the line segments $\overline{AB}$, $\overline{BC}$, $\overline{CD}$, $\overline{DE}$, and $\overline{EA}$ form an arithmetic sequence, although not necessarily in this order. What is the middle term of the arithmetic sequence? 


 
                    \begin{figure}[H]
                        \centering
                        \includegraphics[width=0.5\textwidth]{Problem/AMC_10_10A_2005_17_Problem_Figure_1.png}
                    \end{figure}
                    


 $\textbf{(A) } 9\qquad \textbf{(B) } 10\qquad \textbf{(C) } 11\qquad \textbf{(D) } 12\qquad \textbf{(E) } 13$


 
 \textbf{Answer: }D



\textbf{Solution 1 (assuming a solution exists)}

Each of $A$, $B$, $C$, $D$, and $E$ forms part of exactly $2$ sums along line segments (e.g. $A$ forms part of the sums of $\overline{AB}$ and $\overline{EA}$). Thus, the sum of all these line segments - i.e. the sum of the $5$ terms of the arithmetic sequence - is exactly twice the sum of $A$, $B$, $C$, $D$, and $E$, and since those numbers are $3$, $5$, $6$, $7$, and $9$ in some order, this reduces to $2 \cdot (3+5+6+7+9) = 60$.

 Now, since the middle term of an arithmetic sequence with an odd number of terms is the average of all the terms of the sequence, the middle number must be $\frac{60}{5} = \boxed{\textbf{(D) } 12}$.

 



\textbf{Solution 2 (assuming a solution exists; quick)}

We observe that the average of $3$, $5$, $6$, $7$, and $9$ is $\frac{3+5+6+7+9}{5} = \frac{30}{5} = 6$. Accordingly, setting all $5$ of $A$, $B$, $C$, $D$, and $E$ to $6$ will give the same sum of sums along the line segments - i.e. the sum of the $5$ terms of the arithmetic sequence (as in Solution 1) - as if we had constructed an actual solution. In particular, this results in all the terms of the arithmetic sequence being $6+6 = 12$, and again as in Solution 1, the middle term must be the average of the $5$ terms, so it is necessarily $\boxed{\textbf{(D) } 12}$.

 



\textbf{Solution 3 (rigorous)}

As the numbers $A$, $B$, $C$, $D$, and $E$ are $3$, $5$, $6$, $7$, and $9$ in some order, we deduce that the smallest term of the arithmetic sequence must be at least $3 + 5 = 8$, while the largest term must be at most $7 + 9 = 16$. 

 Thus, as this sequence has $5$ terms, its common difference $d$ is at most $\frac{16-8}{5-1} = 2$.

 But as $d$ is positive (taking the arithmetic sequence in increasing order without loss of generality) and an integer (since it is a difference of sums of integers), while exactly $2$ of the sums must be odd (since exactly $4$ of the numbers $A$, $B$, $C$, $D$, and $E$ are odd, while the $5$th number is even), it follows that $d$ must be odd. Therefore, the only possibility is $d = 1$.

 Next, again noting that the sequence has exactly $2$ odd terms, and thus $5-2 = 3 > 2$ even terms, we observe that the middle term must have the majority parity, i.e. must be even. This means the $2$ terms adjacent to the middle term must be odd, and therefore must be $(6+a)$ and $(6+b)$ with $a,b \in \{3,5,7,9\}$. Moreover, $(6+a)-(6-b) = a-b$ is exactly twice the common difference, i.e. $2d = 2 \cdot 1 = 2$.

 The only other observation we need is that $a$ and $b$ can't be the smallest, or largest, $2$ odd numbers, because that would make it impossible to construct the smallest, or (respectively) largest, sum from $1$ of the remaining $2$ even numbers and $1$ of the odd numbers. Together with $a-b = 2$ from above, this implies that $a = 5$ and $b = 7$, so finally, the middle term of the sequence must be $(6+5) + 1 = (6+7) - 1 = \boxed{\textbf{(D) } 12}$.

 (It is now easy to see that one possible solution is $A = 3$, $B = 9$, $C = 5$, $D = 6$, and $E = 7$; this gives the edge sums as $12$, $14$, $11$, $13$, and $10$, which indeed form the arithmetic sequence $10,11,12,13,14$.)

 


	\item (\textit{AMC 10 10A 2005}) Team $A$ and team $B$ play a series. The first team to win three games wins the series. Each team is equally likely to win each game, there are no ties, and the outcomes of the individual games are independent. If team $B$ wins the second game and team $A$ wins the series, what is the probability that team $B$ wins the first game? 


 $\textbf{(A) } \frac{1}{5}\qquad \textbf{(B) } \frac{1}{4}\qquad \textbf{(C) } \frac{1}{3}\qquad \textbf{(D) } \frac{1}{2}\qquad \textbf{(E) } \frac{2}{3}$


 
 \textbf{Answer: }A



\textbf{Solution}

There are at most $5$ games played. 

 If team $B$ won the first two games, team $A$ would need to win the next three games. So the only possible order of wins is $BBAAA$. 

 If team $A$ won the first game, and team $B$ won the second game, the possible orders of wins are $ABBAA$, $ABABA$, and $ABAAX$, where $X$ denotes that the $5$th game wasn't played. 

 There is $1$ possibility where team $B$ wins the first game, and $4$ total possibilities when team $A$ wins the series and team $B$ wins the second game. Note that the fourth possibility, $ABAAX$, occurs twice as often as the others because it is dependent on the outcome of only $4$ games rather than $5$, so we put $1$ over $1 \cdot 2 + (4-1) = 5$ total possibilities. The desired probability is then $\boxed{\textbf{(A) } \frac{1}{5}}$.

 \textbf{Note} The original problem was poorly worded, since the problem directly stated that the answer is $\boxed{1/2}$ (the probability of team $B$ simply winning the first game, as opposed to the desired conditional probability of this event \textit{given that} both team $B$ wins the second game and team $A$ wins the series). 

 Although the conditional probability is implied by the part of the part of the question saying "If team $B$ wins...", the problem can be better understood as: "What fraction of possible sets of game outcomes have $B$ winning the first game?" or "Given the observed results, what is the conditional probability that $B$ won the first game?"

 


	\item (\textit{AMC 10 10A 2005}) Three one-inch squares are placed with their bases on a line. The center square is lifted out and rotated $45^{\circ}$, as shown. Then it is centered and lowered into its original location until it touches both of the adjoining squares. How many inches is the point $B$ from the line on which the bases of the original squares were placed?


 
                    \begin{figure}[H]
                        \centering
                        \includegraphics[width=0.5\textwidth]{Problem/AMC_10_10A_2005_19_Problem_Figure_1.png}
                    \end{figure}
                    


 $\textbf{(A) } 1\qquad\textbf{(B) } \sqrt{2}\qquad\textbf{(C) } \frac{3}{2}\qquad\textbf{(D) } \sqrt{2}+\frac{1}{2}\qquad\textbf{(E) } 2$


 
 \textbf{Answer: }D



\textbf{Solution}

The rotated middle square is lowered until it touches both the adjoining squares, so since the horizontal distance between those squares is $1$ inch, the middle square will stop being lowered once the length $DE$ in the diagram below is $1$ inch. Now, since $DC$ and $EC$ respectively pass through the vertices $D$ and $E$ of the horizontal middle square, and intersect at right angles, they must be the diagonals of the horizontal middle square, so $\triangle DEC$ is a $45^{\circ}-45^{\circ}-90^{\circ}$ triangle.

 
                    \begin{figure}[H]
                        \centering
                        \includegraphics[width=0.5\textwidth]{Solution/AMC_10_10A_2005_19_Solution_Figure_1.png}
                    \end{figure}
                    

 It follows that when $DE = 1$, we have $DC = EC = \frac{1}{\sqrt{2}} = \frac{\sqrt{2}}{2}$, and as the middle square was rotated by exactly $45^{\circ}$ from its original horizontal position, the diagonal $BC$ is now vertical, giving $BC \perp DE$. This means $\triangle DFC$ and $\triangle EFC$ are also $45^{\circ}-45^{\circ}-90^{\circ}$ triangles, and hence $FC = \frac{\left(\frac{\sqrt{2}}{2}\right)}{\sqrt{2}} = \frac{1}{2}$.

 Lastly, since $BC$ is the diagonal of a unit square, its length is $\sqrt{1^2+1^2} = \sqrt{2}$ (by Pythagoras' theorem), so we deduce that the distance from $B$ to the bottom horizontal line is 

 
 \[BC-FC+1 = \sqrt{2}-\frac{1}{2}+1=\boxed{\textbf{(D) }\sqrt{2}+\dfrac{1}{2}}.\]

 (An alternative for this last step is to compute the distance from $C$ to the bottom horizontal line as $1-FC = 1-\frac{1}{2} = \frac{1}{2}$, so that the distance from $B$ to the bottom horizontal line is, once again, $BC+\frac{1}{2} = \boxed{\textbf{(D) } \sqrt{2}+\frac{1}{2}}$.)

 


	\item (\textit{AMC 10 10A 2005}) An equiangular octagon has four sides of length $1$ and four sides of length $\sqrt{2}/2$, arranged so that no two consecutive sides have the same length.  What is the area of the octagon?


 $\textbf{(A) } \frac{7}{2}\qquad \textbf{(B) } \frac{7\sqrt{2}}{2}\qquad \textbf{(C) } \frac{5+4\sqrt{2}}{2}\qquad \textbf{(D) } \frac{4+5\sqrt{2}}{2}\qquad \textbf{(E) } 7$


 
 \textbf{Answer: }A



\textbf{Solution 1}

The sum of the octagon's angles is $180\cdot(8-2)^{\circ} = 1080^{\circ}$, so since it is equiangular, each angle is $\frac{1080^{\circ}}{8} = 135^{\circ}$, i.e. the same as in a regular octagon. This means that this octagon, just like with a regular octagon, can be divided into squares (whose diagonals form $45^{\circ}$ angles) and right triangles.

 In particular, as shown in the diagram below, the area of this octagon can be divided up into $5$ squares of side length $\frac{\sqrt{2}}{2}$ and $4$ right triangles, each of which has half the area of each of the squares (since each square has diagonal $\frac{\sqrt{2}}{2} \cdot \sqrt{2} = 1$, so each right triangle is congruent by side-side-side to one of the triangles formed by slicing a square along one of its diagonals).

 Therefore, the area of the octagon is equal to the area of $5 + 4 \cdot \frac{1}{2} = 7$ of the squares, and the area of each square is simply $\left(\frac{\sqrt{2}}{2}\right)^2 = \frac{1}{2}$, so the answer is is $7 \cdot \frac{1}{2} = \boxed{\textbf{(A) } \frac{7}{2}}$.

 
                    \begin{figure}[H]
                        \centering
                        \includegraphics[width=0.5\textwidth]{Solution/AMC_10_10A_2005_20_Solution_Figure_1.png}
                    \end{figure}
                    

 



\textbf{Solution 2}

Using the same diagram as in Solution 1, we can extend the sides of length $1$ to form a bounding square around the octagon, with the region inside the square but outside the octagon consisting of $4$ right triangles. (We could also extend the sides of length $\frac{\sqrt{2}}{2}$ to give a similar bounding square, but this would make the computations slightly harder.)

 As in Solution 1, since all the interior angles of the octagon are $135^{\circ}$, the right triangles are all $45^{\circ}-45^{\circ}-90^{\circ}$ triangles with hypotenuse $\frac{\sqrt{2}}{2}$. Thus the length of their other legs is $\frac{\left(\frac{\sqrt{2}}{2}\right)}{\sqrt{2}} = \frac{1}{2}$, so the bounding square has total side length $1 + 2 \cdot \frac{1}{2} = 2$, and hence area $2^2 = 4$. Each of the right triangles has area $\frac{1}{2} \cdot \left(\frac{1}{2}\right)^2 = \frac{1}{8}$, so we deduce that the area of the octagon is $4 - 4 \cdot \frac{1}{8} = \boxed{\textbf{(A) } \frac{7}{2}}$.

 edited by mobius247

 



\textbf{Solution 3}

As shown in the diagram below, we can also divide the octagon into squares and right triangles using horizontal and vertical lines, rather than the lines at a $45^{\circ}$ angle that were used in Solutions 1 and 2.

 
                    \begin{figure}[H]
                        \centering
                        \includegraphics[width=0.5\textwidth]{Solution/AMC_10_10A_2005_20_Solution_Figure_2.png}
                    \end{figure}
                    

 (Diagram made using Geogebra)

 For the same reasons as in Solutions 1 and 2, the octagon's $135^{\circ}$ interior angles mean that the right triangles are $45^{\circ}-45^{\circ}-90^{\circ}$ triangles with hypotenuse $\frac{\sqrt{2}}{2}$, so their other legs have length $\frac{\left(\frac{\sqrt{2}}{2}\right)}{\sqrt{2}} = \frac{1}{2}$. Accordingly, their areas are each $\frac{1}{2} \cdot \left(\frac{1}{2}\right)^2 = \frac{1}{8}$, while by symmetry, each of the squares also has side length $\frac{1}{2}$, and thus area $\left(\frac{1}{2}\right)^2 = \frac{1}{4}$.

 Hence, as the octagon consists of $12$ squares and $4$ right triangles, its total area is $12\cdot\frac{1}{4}+4\cdot\frac{1}{8} = \boxed{\textbf{(A) } \frac{7}{2}}$.

 



\textbf{Solution 4}

We notice that we can draw four $45-45-90$ right triangles(each with side length $\frac{1}{2}$) and make the octagon a square with side length $2$. 

 
                    \begin{figure}[H]
                        \centering
                        \includegraphics[width=0.5\textwidth]{Solution/AMC_10_10A_2005_20_Solution_Figure_3.png}
                    \end{figure}
                    

 The area of the resulting square is $2^2 = 4$. Additionally, the area of each extra triangle is $\frac{1}{2} \cdot \frac{1}{2} \cdot \frac{1}{2} = \frac{1}{8}$, and so four of them would be $4 \cdot \frac{1}{8} = \frac{4}{8} = \frac{1}{2}$. Therefore, the area of the octagon(which can be gotten by subtracting the area of the triangles from the area of the square) is $4 - \frac{1}{2} = \boxed{\textbf{(A) } \frac{7}{2}}$.

 


	\item (\textit{AMC 10 10A 2005}) For how many positive integers $n$ does $1+2+\dotsb+n$ evenly divide $6n$? 


 $\textbf{(A) } 3\qquad \textbf{(B) } 5\qquad \textbf{(C) } 7\qquad \textbf{(D) } 9\qquad \textbf{(E) } 11$


 
 \textbf{Answer: }B



\textbf{Solution}

By a standard result, $1+2+\dotsb+n = \frac{n(n+1)}{2}$, so this will evenly divide $6n$ precisely if $\frac{6n}{\left(\frac{n(n+1)}{2}\right)} = \frac{12}{n+1}$ is an integer, or equivalently, $(n+1)$ is a factor of $12$. 

 The positive factors of $12$ are $1$, $2$, $3$, $4$, $6$, and $12$ (i.e. $6$ factors in total), but as $n$ must be a positive integer, we must have $n+1 \geq 1+1 = 2$. Therefore, the factor $1$ cannot be used, but all the others can, giving $6-1 = \boxed{\textbf{(B) } 5}$ possible values of $n$ (namely, $n$ can be $2-1 = 1$, $3-1 = 2$, $4-1 = 3$, $6-1 = 5$, or $12-1 = 11$).

 


	\item (\textit{AMC 10 10A 2005}) Let $S$ be the set of the $2005$ smallest positive multiples of $4$, and let $T$ be the set of the $2005$ smallest positive multiples of $6$. How many elements are common to $S$ and $T$?


 $\textbf{(A) } 166\qquad \textbf{(B) } 333\qquad \textbf{(C) } 500\qquad \textbf{(D) } 668\qquad \textbf{(E) } 1001$


 
 \textbf{Answer: }D



\textbf{Solution 1}

Since the least common multiple of $4$ and $6$ is $12$, the elements that are common to $S$ and $T$ are all multiples of $12$. Moreover, as the largest element of $S$ is $4 \cdot 2005$, while that of $T$ is $6 \cdot 2005$, which is larger, several of the multiples of $12$ that are in $T$ will not be in $S$, whereas all the multiples of $12$ that are in $S$ will be in $T$.

 Thus we only need to find the number of multiples of $12$ that are in $S$, and so we observe that as $4 \cdot 3 = 12$, these multiples of $12$ are precisely every $3$rd element of $S$. It follows that there are $\left\lfloor\frac{2005}{3}\right \rfloor = \boxed{\textbf{(D) } 668}$ such elements.

 



\textbf{Solution 2}

As in Solution $1$, we find that the elements common to $S$ and $T$ are precisely the multiples of $12$. As $S$ has exactly $2005$ elements, these must range from $4 \cdot 1 = 4$ to $4 \cdot 2005 = 8020$, and similarly the elements of $T$ range from $6 \cdot 1 = 6$ to $6 \cdot 2005 = 12030$. This means any element $n \in S \cap T$ must satisfy both $4 \leq n \leq 8020$ and $6 \leq n \leq 12030$, which reduces to simply $6 \leq n \leq 8020$.

 Accordingly, as $12 \cdot 0 = 0 < 6 < 12 \cdot 1 = 12$ and $12 \cdot 668 = 8016 < 8020 < 12 \cdot 669 = 8028$, the multiples of $12$ in the required interval range from $12 \cdot 1$ to $12 \cdot 668$, so there are precisely $\boxed{\textbf{(D) } 668}$ of them.

 


	\item (\textit{AMC 10 10A 2005}) Let $\overline{AB}$ be a diameter of a circle and $C$ be a point on $\overline{AB}$ with $2 \cdot AC = BC$. Let $D$ and $E$ be points on the circle such that $\overline{DC} \perp \overline{AB}$ and $\overline{DE}$ is a second diameter. What is the ratio of the area of $\triangle DCE$ to the area of $\triangle ABD$?


 
                    \begin{figure}[H]
                        \centering
                        \includegraphics[width=0.5\textwidth]{Problem/AMC_10_10A_2005_23_Problem_Figure_1.png}
                    \end{figure}
                    


 $\textbf{(A) } \frac{1}{6} \qquad \textbf{(B) } \frac{1}{4} \qquad \textbf{(C) } \frac{1}{3} \qquad \textbf{(D) } \frac{1}{2} \qquad \textbf{(E) } \frac{2}{3}$


 
 \textbf{Answer: }C



\textbf{Solution 1}


                    \begin{figure}[H]
                        \centering
                        \includegraphics[width=0.5\textwidth]{Solution/AMC_10_10A_2005_23_Solution_Figure_1.png}
                    \end{figure}
                    

 WLOG, let us assume that the diameter is of length $1$. 

 The length of $AC$ is $\frac{1}{3}$ and $CO$ is $\frac{1}{2}-\frac{1}{3} = \frac{1}{6}$. 

 $OD$ is the radius of the circle, which is $\frac{1}{2}$, so using the \href{https://artofproblemsolving.com/wiki/index.php?title=Pythagorean\_Theorem}{Pythagorean Theorem} the height $CD$ of $\triangle DCO$ is $\sqrt{\left(\frac{1}{2}\right)^2-\left(\frac{1}{6}\right)^2} = \frac{\sqrt{2}}{3}$. This is also the height of the $\triangle ABD$.

 The area of $\triangle DCO$ is $\frac{1}{2}\cdot\frac{1}{6}\cdot\frac{\sqrt{2}}{3}$ = $\frac{\sqrt{2}}{36}$.

 The height of $\triangle DCE$ can be found using the area of $\triangle DCO$ and $DO$ as base.

 Hence, the height of $\triangle DCE$ is $\dfrac{\dfrac{\sqrt{2}}{36}}{\dfrac{1}{2}\cdot\dfrac{1}{2}}$ = $\dfrac{\sqrt{2}}{9}$. 

 The diameter is the base for both the triangles $\triangle DCE$ and $\triangle ABD$,

 Hence, the ratio of the area of $\triangle DCE$ to the area of $\triangle ABD$ is$\dfrac{\dfrac{\sqrt{2}}{9}}{\dfrac{\sqrt{2}}{3}}$ = $\boxed{\textbf{(C) }\frac{1}{3}}$

 



\textbf{Solution 2}

Since $\triangle DCE$ and $\triangle ABD$ share a base, the ratio of their areas is the ratio of their altitudes. Draw the altitude from $C$ to $DE$. 

 
                    \begin{figure}[H]
                        \centering
                        \includegraphics[width=0.5\textwidth]{Solution/AMC_10_10A_2005_23_Solution_Figure_2.png}
                    \end{figure}
                    $OD=r, OC=\frac{1}{3}r$.

 Since $m\angle DCO=m\angle DFC=90^\circ$, then $\triangle DCO\sim \triangle DFC$. So the ratio of the two altitudes is $\frac{CF}{DC}=\frac{OC}{DO}=\boxed{\textbf{(C) }\frac{1}{3}}$

 



\textbf{Solution 3}

Let the center of the circle be point $O$;Without loss of generality, assume $AC=2$, so $CB=4$ and the diameter and radius are $6$ and $3$, respectively. Therefore, $CO=1$, and $DO=3$.The area of $\triangle DCE$ can be expressed as $\frac{1}{2}(CD)(6)\text{sin }(CDE).$ $\frac{1}{2}(CD)(6)$ happens to be the area of $\triangle ABD$. Furthermore, $\text{sin } CDE = \frac{CO}{DO},$ or $\frac{1}{3}.$ Therefore, the ratio is $\boxed{\textbf{(C) }\frac{1}{3}}.$

 



\textbf{Solution 4}


                    \begin{figure}[H]
                        \centering
                        \includegraphics[width=0.5\textwidth]{Solution/AMC_10_10A_2005_23_Solution_Figure_3.png}
                    \end{figure}
                    

 


 Let the point G be the reflection of point $D$ across $\overline{AB}$. (Point G is on the circle).

 
Let $AC=x$, then $BC=2x$. The diameter is $3x$. To find $DC$, there are two ways (presented here):

 1. Since $\overline{AB}$ is the diameter, $CD=CG$. Using power of points,
 \[AC\cdot BC=x\cdot2x=2x^{2}=CD^{2} \longrightarrow CD=x\sqrt{2}\]2. Use the geometric mean theorem,
 \[AC\cdot BC=x\cdot2x=2x^{2}=CD^{2} \longrightarrow CD=x\sqrt{2}\](These are the same equations but obtained through different formulae)

 
Therefore $DG=2x\sqrt{2}$. Since $\overline{DE}$ is a diameter, $\triangle DGE$ is right. By the Pythagorean theorem,
 \[DE^{2}=GD^{2}+GE^{2} \longrightarrow \left(3x\right)^{2}=\left(2x\sqrt{2}\right)^{2}+GE^{2}\]
 \[9x^{2}=8x^{2}+GE^{2} \longrightarrow GE^{2}=x^{2} \longrightarrow GE=x\]

 
As established before, $\angle DGE$ is right (if you are unsure, look up "inscribed angle theorem", this is a special case of the theorem where the central angle measures $180^{\circ}$) so $GE=x$ is the altitude of $\triangle DCE$, and $DC=x\sqrt{2}$ is the base. Therefore 
 \[\left[DCE\right]=\frac{1}{2}\cdot DC\cdot GE=\frac{1}{2}\cdot x\sqrt{2}\cdot x=\frac{x^{2}\sqrt{2}}{2}\]

 
$AB=3x$ is the base of $\triangle ABD$ and $CD=x\sqrt{2}$ is the height.
 \[\left[ABD\right]=\frac{1}{2}\cdot3x\cdot x\sqrt{2}=\frac{3x^{2}\sqrt{2}}{2}\]

 
The required ratio is 
 \[\frac{\left[DCE\right]}{\left[ABD\right]}=\frac{\frac{x^{2}\sqrt{2}}{2}}{\frac{3x^{2}\sqrt{2}}{2}}=\frac{x^{2}\sqrt{2}}{2}\cdot\frac{2}{3x^{2}\sqrt{2}}=\frac{x^{2}\sqrt{2}}{3x^{2}\sqrt{2}}=\frac{1}{3}\]The answer is $\boxed{\textbf{(C) } \frac{1}{3}}$.

 


 



\textbf{Solution 5}

Assume the diameter is $3$.

 $AC = 1$

 Get the height $CD = \sqrt{(AC)(BC)} = \sqrt2$ via power of a point.

 $CO = AO - AC = 1/2$.

 By altitude of right triangle $\triangle CDO$: Altitude from $C$ to $DE$ is same as altitude from $C$ to $DO$ is $\frac{(CO)(CD)}{DO} = \frac{(1/2)(\sqrt2)}{\frac 3 2}$.

 $\triangle DCE$ and $\triangle ABD$ have the same (diameter) hypotenuse length, so their area ratio is their altitude ratio is$\frac {\frac{ (1/2)(\sqrt2) } {\frac 3 2}} {\sqrt2} = \boxed{1/3}$.

 



\textbf{Solution 6}

Think about $DE$ and $AB$ intersecting at $X$. Now we realize that $XA = XB = XE = XD$ (this distance is just the radius of the circle), so that means that $X$ is the center.

 So this can help us conclude that $\triangle ABE$ and $\triangle ADB$ are right triangles, and since it’s a rectangle setup, by Side-Angle-Side they are similar:$\triangle ABE \sim \triangle ADB.$

 So because of this, the altitudes of $\triangle DCE$, $\triangle ABE$, and $\triangle ADB$ are the same (with $CD$ as the altitude). Hence we need to find the ratio$\frac{2CX}{AB}.$

 We have 
 \[AC:CB = 1:2 \Rightarrow \quad AX:XB = 1:1 \Rightarrow AC:CX = 2:1.\]

 Since $AC:CB = 1:2$, it follows that
 \[BX:CX = 1:4 \Rightarrow CX:AB = 1:6.\]

 Therefore,
 \[\frac{2CX}{AB} = 2 \cdot \frac{1}{6} = \boxed{\frac{1}{3}}\]

 


	\item (\textit{AMC 10 10A 2005}) For each positive integer $m > 1$, let $P(m)$ denote the greatest prime factor of $m$. For how many positive integers $n$ is it true that both $P(n) = \sqrt{n}$ and $P(n+48) = \sqrt{n+48}$?


 $\textbf{(A) } 0\qquad \textbf{(B) } 1\qquad \textbf{(C) } 3\qquad \textbf{(D) } 4\qquad \textbf{(E) } 5$


 
 \textbf{Answer: }B



\textbf{Solution 1}

If $P(n) = \sqrt{n}$, then $n = p_{1}^{2}$, where $p_{1}$ is a \href{https://artofproblemsolving.com/wiki/index.php?title=Prime\_number}{prime number}.

 If $P(n+48) = \sqrt{n+48}$, then $n + 48$ is a square, but we know that n is $p_{1}^{2}$.

 
This fact reveals that we just have to check for squares of primes, add $48$ and look whether the root is a prime number.We can easily see that the difference between two consecutive square after $576$ is greater than or equal to $49$,Hence we have to consider only the prime numbers till $23$.

 Squaring prime numbers below $23$ including $23$ we get the following list.

 $4 , 9 , 25 , 49 , 121, 169 , 289 , 361 , 529$

 But adding $48$ to a number ending with $9$ will result in a number ending with $7$, but we know that a perfect square does not end in $7$, so we can eliminate those cases to get the new list.

 $4 , 25 , 121 , 361$

 Adding $48$, we get $121$ as the only possible solution.Hence the answer is $\boxed{\textbf{(B) }1}$.

 edited by mobius247, ChristianYoo

 \textbf{Note: Solution 1} Since all primes greater than $2$ are odd, we know that the difference between the squares of any two consecutive primes greater than $2$ is at least $(p+2)^2-p^2=4p+4$, where p is the smaller of the consecutive primes. For $p>11$, $4p+4>48$. This means that the difference between the squares of any two consecutive primes both greater than $11$ is greater than $48$, so $n$ and $n+48$ can't both be the squares of primes if $n=p^2$ and $p>11$. So, we only need to check $n=2^2, 3^2, 5^2, 7^2,$ and $11^2$.

 


	\item (\textit{AMC 10 10A 2005}) In $\triangle ABC$ we have $AB = 25$, $BC = 39$, and $AC = 42$. Points $D$ and $E$ are on $\overline{AB}$ and $\overline{AC}$ respectively, with $AD = 19$ and $AE = 14$. What is the ratio of the area of triangle $ADE$ to the area of the quadrilateral $BCED$?


 $\textbf{(A) } \frac{266}{1521}\qquad \textbf{(B) } \frac{19}{75}\qquad \textbf{(C) } \frac{1}{3}\qquad \textbf{(D) } \frac{19}{56}\qquad \textbf{(E) } 1$


 
 \textbf{Answer: }D



\textbf{Solution 1}

We have
 \[\frac{[ADE]}{[ABC]} = \frac{AD}{AB} \cdot \frac{AE}{AC} = \frac{19}{25} \cdot \frac{14}{42} = \frac{19}{75}.\]

 (Area of a triangle is half times base times height, so the area ratio of triangles, that have a common vertex (height) and bases on a common line, is the base length ratio. This is applied twice, using different pairs of bases, and corresponding altitudes for height.).

 
                    \begin{figure}[H]
                        \centering
                        \includegraphics[width=0.5\textwidth]{Solution/AMC_10_10A_2005_25_Solution_Figure_1.png}
                    \end{figure}
                    

 $[BCED] = [ABC] - [ADE]$, so
 \begin{align*} \frac{[ADE]}{[BCED]} &= \frac{[ADE]}{[ABC] - [ADE]} \\ &= \frac{1}{\frac{ABC}{ADE} - 1} \\ &= \frac{1}{\frac{75}{19} - 1} \\ &= \boxed{\textbf{(D) }\frac{19}{56}}. \end{align*}


 


 Note: If it is hard to understand why 
 \[\frac{[ADE]}{[ABC]} = \frac{AD}{AB} \cdot \frac{AE}{AC}\], you can use the fact that the area of a triangle equals $\frac{1}{2} \cdot ab \cdot \sin(C)$. If angle $DAE = Z$, we have that 
 \[\frac{[ADE]}{[ABC]} = \frac{\frac{1}{2} \cdot 19 \cdot 14 \cdot \sin(Z)}{\frac{1}{2} \cdot 25 \cdot 42 \cdot \sin(Z)} = \frac{19 \cdot 14}{25 \cdot 42} = \frac{ab}{cd}\].

 


\end{enumerate}